\documentclass[10pt,a4paper]{amsart}
\usepackage[margin=19mm]{geometry}
\usepackage{amsmath,amssymb,mathtools}
\usepackage[scr=rsfs,cal=euler]{mathalfa}
\usepackage{xfrac}
\usepackage[unicode,hidelinks,bookmarks=false]{hyperref}
\newcommand{\ie}{i.e.\ }
\newcommand{\cf}{cf.\ }
\newcommand{\resp}{respectively\ }
\newcommand{\Subsection}{\subsection}
\providecommand{\nobreakdash}{\nobreak\hbox{-}}
\setlength{\emergencystretch}{2em}
\multlinegap0pt
\usepackage{bm}
\usepackage{bbm}
\usepackage{dsfont}
\usepackage{xspace}
\usepackage{cancel}
\usepackage{mathtools}
\usepackage{amsthm}
\makeatletter
\@ifundefined{theorem}{\newtheorem{theorem}{Theorem}}{}
\makeatother
\def\ce{\mathcal{E}}
\def\ci{\mathcal{I}}
\newcommand{\D}{\mathcal{D}}
\newcommand{\R}{\mathbb{R}}
\title[Heat kernel bounds for the fractional Laplacian with Hardy p]{Heat kernel bounds for the fractional Laplacian with Hardy potential in angular momentum channels}
\author{Krzysztof Bogdan, Konstantin Merz}
\date{Source: arXiv:2506.08115v1 (CC BY 4.0)}
\begin{document}
\maketitle

\noindent\emph{Source and licence.} Heat kernel bounds for the fractional Laplacian with Hardy potential in angular momentum channels. Version arXiv:2506.08115v1 (CC BY 4.0), distributed under the Creative Commons Attribution 4.0 International licence, as recorded in the arXiv licence metadata for this exact version and shown on the abstract page https://arxiv.org/abs/2506.08115v1. Full rights notice: licenses/arxiv-2506.08115-CC-BY-4.0.txt.\par\smallskip

\noindent\textit{Editorial scope.} Counted unit: the Theorem whose source label is \texttt{jakprop318} in the article (source0.tex lines 1723--1730, source label \texttt{jakprop318}), delivered together with the article's own complete original proof of it (source0.tex lines 1732--1772, one \texttt{proof} environment of 41 lines). No mathematics is written, repaired or re-derived: statement and proof are the article's own text line for line. Required context retained uncounted: propertiesschrodheatkernel (lines 777--818); eq:defnuell4 (lines 866--870); eq:nuellrepulsive (lines 1702--1706). Every definition, display and label of those lines is retained unchanged. Omitted from this package and not redistributed: 1--776, 819--865, 871--1701, 1707--1722, 1731--1731, 1773--2249. Nothing inside a retained excerpt is omitted or altered: every retained body is delivered exactly as the article's lines are written, so no omission receipt of any kind accompanies this package. Citations: the retained excerpts carry no citation command at all, so no bibliography entry of the article is transcribed and no reference is redistributed. Reference adaptation: none was needed and none was made; every reference inside the delivered unit resolves inside it, so no unresolved reference or citation is produced. Printed slips kept literal: no printed word, symbol, hypothesis or number of the counted statement or of its proof was altered. Layout and class: the article's own class is \texttt{amsart} with packages including amsmath, amsfonts, amsthm, amssymb, amsxtra, hyperref, bbm, bm, graphicx, float, caption; the wrapper is the lane's pinned \texttt{amsart} editorial wrapper at 10pt on A4, which restates the theorem-like environments the retained text needs and adds the article's own macro definitions that the retained text needs (\texttt{\textbackslash D}, \texttt{\textbackslash R}, \texttt{\textbackslash ce}, \texttt{\textbackslash ci}), with extra wrapper packages (bm, bbm, dsfont, xspace, cancel, mathtools). The delivered numbering is the unit's own. Assets: no retained span contains a figure environment, a graphics-inclusion command, a PDF-page inclusion, a table or a TikZ picture; no image file is copied into this package, the package declares no asset dependency and no image file is redistributed with it. Licence: version arXiv:2506.08115v1 of this article is distributed under the Creative Commons Attribution 4.0 International licence, as recorded in the arXiv licence metadata for this exact version and shown on the abstract page https://arxiv.org/abs/2506.08115v1. A full rights notice is delivered at licenses/arxiv-2506.08115-CC-BY-4.0.txt. Characters: every retained excerpt is pure ASCII, so the delivered engine, XeLaTeX with its default Latin Modern text font, reports no missing character.
\begin{theorem}
  \label{propertiesschrodheatkernel}
  Let $\zeta\in(-1/2,\infty)$, $\alpha\in(0,2)\cap(0,2\zeta+1)$, and $\eta\in(-\alpha,\frac{2\zeta+1-\alpha}{2}]$, $r,s,t,t'>0$. Then the following statements hold.\\
  \textup{(1)}
  We have the Chapman--Kolmogorov equation
  \begin{align}
    \int_0^\infty dz\, z^{2\zeta} p_{\zeta,\eta}^{(\alpha)}(t,r,z) \, p_{\zeta,\eta}^{(\alpha)}(t',z,s)
    & = p_{\zeta,\eta}^{(\alpha)}(t+t',r,s).
  \end{align}
  \textup{(2)}
  We have the Duhamel formulae
  \begin{align}
    \label{eq:duhamelclassictransformed}
    \begin{split}
      p_{\zeta,\eta}^{(\alpha)}(t,r,s)
      & = p_\zeta^{(\alpha)}(t,r,s) + \int_0^t d\tau \int_0^\infty dz\, z^{2\zeta} p_\zeta^{(\alpha)}(\tau,r,z) q(z)\, p_{\zeta,\eta}^{(\alpha)}(t-\tau,z,s) \\
      & = p_\zeta^{(\alpha)}(\tau,r,s) + \int_0^t d\tau \int_0^\infty dz\, z^{2\zeta} p_{\zeta,\eta}^{(\alpha)}(\tau,r,z) q(z)\, p_\zeta^{(\alpha)}(t-\tau,z,s).
    \end{split}
  \end{align}
  \textup{(3)}
  We have the scaling relation
  \begin{align}
    \label{eq:scalingalphahardy}
    p_{\zeta,\eta}^{(\alpha)}(t,r,s)
    = t^{-\frac{2\zeta+1}{\alpha}} p_{\zeta,\eta}^{(\alpha)}\left(1,\frac{r}{t^{1/\alpha}},\frac{s}{t^{1/\alpha}}\right).
  \end{align}
  \textup{(4)}
  The function $h(r)=r^{-\eta}$ is invariant under $p_{\zeta,\eta}^{(\alpha)}$, i.e.,
  \begin{align}
    \label{eq:supermediantransformedalphaintro}
    \begin{split}
      % \int_0^\infty ds\, s^{2\zeta} p_\zeta^{(\alpha)}(t,r,s) h(s)
      % & \leq h(r), \\
      \int_0^\infty ds\, s^{2\zeta} p_{\zeta,\eta}^{(\alpha)}(t,r,s) h(s) = h(r), \quad t>0.
    \end{split}
  \end{align}
  \textup{(5)}
  For $\eta>0$, $p_\zeta^{(\alpha)}(t,r,s)\leq p_{\zeta,\eta}^{(\alpha)}(t,r,s)$, and for $\eta<0$, $p_\zeta^{(\alpha)}(t,r,s)\geq p_{\zeta,\eta}^{(\alpha)}(t,r,s)$.
  \\
  \textup{(6)}
  $\{p_{\zeta,\eta}^{(\alpha)}(t,\cdot,\cdot)\}_{t\geq0}$ is a strongly continuous semigroup of contractions on $L^2(\R_+,r^{2\zeta}dr)$.
\end{theorem}

\begin{align}
  \label{eq:defnuell4}
  \frac{p_\zeta^{(\alpha)}(t,r,s)}{t}
  \lesssim_{\zeta,\alpha} \nu_{\zeta}(r,s), \quad r,s,t>0.
\end{align}

  \begin{align}
    \label{eq:nuellrepulsive}
    \lim_{t\to0}\frac{p_{\zeta,\eta}^{(\alpha)}(t,r,s)}{t}
    = \nu_\zeta(r,s),\quad r,s>0,\ r\neq s.
  \end{align}

\begin{theorem}
  \label{jakprop318}
  Let $\zeta\in[0,\infty)$, $\alpha\in(0,2\wedge(2\zeta+1))$, and $\eta\in(-\alpha,0)$.
  Then, $\D(\ce_{\zeta,\eta})=\D(\ci_{\zeta,\eta})$ and
  \begin{align}
    \ce_{\zeta,\eta}[u] = \ci_{\zeta,\eta}[u],\quad u\in\D(\ce_{\zeta,\eta}).
  \end{align}
\end{theorem}

\begin{proof}
  % Recall that $\{p_{\zeta,\eta}^{(\alpha)}\}_{t\geq0}$ is well-defined on $L^2(\R_+,r^{2\zeta}dr)$ by setting $(p_{\zeta,\eta}^{(\alpha)}(t,\cdot,\cdot)u)(0)=0$ for any $u\in L^2(\R_+,r^{2\zeta}dr)$.
  By Theorem~\ref{propertiesschrodheatkernel}(4),
  \begin{align}
    \begin{split}
      & \langle u,u-p_{\zeta,\eta}^{(\alpha)}(t,\cdot,\cdot)u\rangle_{L^2(\R_+,r^{2\zeta}dr)}
      = \int_{\R_+}dr\, r^{2\zeta}\, \overline{u(r)}\left(u(r) - \int_{\R_+}ds\, s^{2\zeta} p_{\zeta,\eta}^{(\alpha)}(t,r,s) u(s)\right) \\
      & \quad = \int_{\R_+}dr\, r^{2\zeta}\, \overline{u(r)}\left(\frac{u(r)}{h(r)} \int_0^\infty ds\, s^{2\zeta} p_{\zeta,\eta}^{(\alpha)}(t,r,s)h(s) - \int_0^\infty ds\, s^{2\zeta} p_{\zeta,\eta}^{(\alpha)}(t,r,s)h(s)\frac{u(s)}{h(s)}\right) \\
      & \quad = \iint_{\R_+\times\R_+}dr\,ds\, (rs)^{2\zeta-\eta}\, \frac{\overline{u(r)}}{h(r)} p_{\zeta,\eta}^{(\alpha)}(t,r,s)\left(\frac{u(r)}{h(r)} - \frac{u(s)}{h(s)}\right).
    \end{split}
  \end{align}
  Hence, by symmetry,
  \begin{align}
    \langle u,u-p_{\zeta,\eta}^{(\alpha)}(t,\cdot,\cdot)u\rangle_{L^2(\R_+,r^{2\zeta}dr)}
    = \frac12 \iint_{\R_+\times\R_+}dr\,ds\, (rs)^{2\zeta-\eta}\, p_{\zeta,\eta}^{(\alpha)}(t,r,s)\left|\frac{u(r)}{h(r)} - \frac{u(s)}{h(s)}\right|^2.
  \end{align}
  Now, for $u\in\D(\ce_{\zeta,\eta})$, by Fatou's lemma and \eqref{eq:nuellrepulsive}, we have
  \begin{align}
    \begin{split}
      \ce_{\zeta,\eta}[u]
      & = \lim_{t\to0}\frac1t\langle u,(1-p_{\zeta,\eta}^{(\alpha)}(t,\cdot,\cdot))u\rangle_{L^2(\R_+,r^{2\zeta}dr)} \\
      & = \lim_{t\to0}\frac12\iint_{\R_+\times\R_+}\frac{p_{\zeta,\eta}^{(\alpha)}(t,r,s)}{t} \left|\frac{u(r)}{h(r)} - \frac{u(s)}{h(s)}\right|^2 (rs)^{2\zeta-\eta}\,dr\,ds \\
      & \geq \frac12\iint\limits_{\R_+\times\R_+}\liminf_{t\to0}\frac{p_{\zeta,\eta}^{(\alpha)}(t,r,s)}{t} \left|\frac{u(r)}{h(r)} - \frac{u(s)}{h(s)}\right|^2 (rs)^{2\zeta-\eta}\,dr\,ds \\
      & = \frac12\iint\limits_{\R_+\times\R_+}dr\,ds\, (rs)^{2\zeta-\eta}\nu_\zeta(r,s) \left|\frac{u(r)}{h(r)} - \frac{u(s)}{h(s)}\right|^2
        = \ci_{\zeta,\eta}[u].
    \end{split}
  \end{align}
  This shows that $\D(\ce_{\zeta,\eta})\subseteq\D(\ci_{\zeta,\eta})$. To prove the reversed inclusion, let $u\in\D(\ci_{\zeta,\eta})$. Then, by $p_{\zeta,\eta}^{(\alpha)}(t,r,s)\leq p_\zeta^{(\alpha)}(t,r,s)\lesssim t\nu_\zeta(r,s)$ for all $r,s,t>0$, \eqref{eq:defnuell4}, by the dominated convergence and \eqref{eq:nuellrepulsive},
  \begin{align}
    \begin{split}
      \ce_{\zeta,\eta}[u]
      & = \lim_{t\to0}\frac1t\langle u,(1-p_{\zeta,\eta}^{(\alpha)}(t,\cdot,\cdot))u\rangle_{L^2(\R_+,r^{2\zeta}dr)} \\
      & = \lim_{t\to0}\frac12 \iint_{\R_+\times\R_+}\frac{p_{\zeta,\eta}^{(\alpha)}(t,r,s)}{t} \left|\frac{u(r)}{h(r)} - \frac{u(s)}{h(s)}\right|^2 h(r)h(s)\, (rs)^{2\zeta}\,dr\,ds \\
      & = \frac12 \iint_{\R_+\times\R_+}dr\,ds\, (rs)^{2\zeta-\eta}\nu_\zeta(r,s) \left|\frac{u(r)}{h(r)} - \frac{u(s)}{h(s)}\right|^2
      = \ci_{\zeta,\eta}[u].
    \end{split}
  \end{align}
  This shows $u\in\D(\ce_{\zeta,\eta})$ and hence $\D(\ci_{\zeta,\eta})\subseteq\D(\ce_{\zeta,\eta})$
  % , and concludes the proof of $\D(\ce_{\zeta,\eta})\subseteq\D(\ci_{\zeta,\eta})$ 
  and $\ce_{\zeta,\eta}[u]=\ci_{\zeta,\eta}[u]$ for all $u\in\D(\ce_{\zeta,\eta})$.
\end{proof}

\end{document}
